\chapter{TINJAUAN PUSTAKA}
\label{bab:tinjauan}

\begin{quote}
\small\itshape
Catatan transkripsi: sebagian besar persamaan bernomor pada bab ini
(\eqref{eq:2-1}--\eqref{eq:2-25}) aslinya disisipkan sebagai objek OLE
\textit{MathType/Equation Editor} di dalam \texttt{BAB II.doc}, bukan sebagai
teks biasa, sehingga tidak dapat diambil langsung sebagai teks oleh
\texttt{antiword}. Persamaan-persamaan tersebut disusun ulang di sini
berdasarkan (a) teks naratif di sekitarnya yang tetap utuh, (b) kode sumber
C++ yang menjadi hasil tugas akhir ini (\texttt{Struktur.hpp},
\texttt{Solver.hpp}, \texttt{Balok.hpp}, \texttt{Kolom.hpp} pada
\texttt{Optimasi Beton/Source/}), dan (c) rujukan standar yang memang
disebut dalam teks (Weaver \& Gere 1980; SK SNI T-15-1991-03; Hulse \&
Mosley 1986). Lihat \texttt{LEGACY\_TRANSCRIPTION\_NOTES.md} untuk detail
metodologinya. Gambar 2-1 dan 2-2 (penampang dan diagram regangan) adalah
citra asli yang tidak dapat dipulihkan.
\end{quote}

\section{Struktur Portal Ruang}

\subsection{Tinjauan Umum}

Langkah pertama dalam mendesain struktur beton bertulang adalah memperoleh
informasi tentang momen, gaya geser, dan gaya aksial yang akan ditahan. Hal
ini biasanya melibatkan perhitungan numerik yang rumit, oleh karena itu
langkah terbaik adalah memprosesnya dengan program analisa struktur yang
telah dibuat dalam komputer (Hulse dan Mosley, 1986, halaman 17).

Dengan perkembangan komputer yang cukup pesat dan harganya relatif murah,
sekarang telah beredar program-program struktur, baik hasil buatan sendiri
maupun perangkat lunak yang diperdagangkan. Dengan perangkat ini suatu
analisis struktur dapat diselesaikan jauh lebih cepat, dan respon sistem
struktur dapat diperoleh. Saat ini analisis tiga dimensi telah umum
digunakan (Wahyudi dan Rahim, 1997, halaman 17).

Dalam menganalisa struktur rangka (\textit{framed structure}) dengan
bantuan komputer dapat digunakan metoda matriks, baik dengan metoda gaya
(\textit{flexibility}) maupun metoda kekakuan (\textit{stiffness}). Metoda
kekakuan atau lebih dikenal sebagai metoda perpindahan lebih sesuai untuk
diprogram dalam komputer, karena setelah model analisis struktur
ditentukan, pertimbangan teknis yang lebih lanjut tidak lagi diperlukan, di
sini terdapat perbedaan dengan metoda gaya, walaupun kedua metoda ini sama
bentuk matematisnya, pada metoda gaya besaran yang tidak diketahui adalah
gaya kelebihan (\textit{redundant}) yang dipilih secara sembarang,
sebaliknya dalam metoda kekakuan, yang tidak diketahui adalah perpindahan
titik kumpul yang tertentu secara otomatis. Jadi, jumlah yang tak diketahui
dalam metoda kekakuan sama dengan derajat ketidaktentuan
(\textit{indeterminacy}) kinematis struktur (Weaver dan Gere, 1980,
halaman 1).

Struktur rangka dibagi atas enam kategori: balok, rangka batang bidang,
rangka batang ruang, portal bidang, balok silang (\textit{grid}) dan portal
ruang, setiap struktur rangka terdiri dari batang-batang yang panjangnya
jauh lebih besar dibandingkan ukuran penampang lintangnya. Titik kumpul
struktur rangka adalah titik pertemuan batang-batang, termasuk tumpuan dan
ujung bebas suatu batang. Tumpuan bisa merupakan jepitan, sendi, atau
tumpuan rol, dalam kasus tertentu sambungan antara batang atau batang
dengan tumpuan bisa bersifat elastis (atau semi-tegar). Beban pada struktur
rangka bisa berupa gaya terpusat, beban tersebar, atau kopel (Weaver dan
Gere, 1980, halaman 2).

\subsection{Metoda Kekakuan}

Metoda kekakuan pertama dikembangkan dari superposisi gaya untuk koordinat
perpindahan bebas, kemudian metoda ini diformalisasi dan diperluas dengan
memakai matriks kesepadanan (\textit{compatibility}) dan konsep kerja maya
(\textit{virtual}). Pada pendekatan kerja maya, matriks kekakuan titik
kumpul (\textit{joint stiffness}) yang lengkap dirakit dengan perkalian
tiga matriks, walaupun versi formal tersebut menarik dan beraturan, tetapi
membutuhkan matriks kesepadanan yang besar dan jarang (\textit{sparse}).
Selain itu, elemen matriks ini tidak mudah dinilai dengan tepat. Baik cara
pembentukan matriks di atas ataupun perakitan dengan proses perkalian tidak
sesuai untuk diprogram pada komputer. Metodologi yang baik adalah menarik
ide dari kedua pendekatan tersebut dan menambahkan beberapa teknik yang
berorientasi pada komputer, sehingga didapatkan suatu cara yang disebut
metoda kekakuan langsung (\textit{direct stiffness method}) yang lebih
mudah untuk diprogram pada komputer (Weaver dan Gere, 1980, halaman 148).

Wahyudi dan Rahim (1997, halaman 17) menyebutkan prosedur umum metoda
kekakuan ataupun metoda elemen hingga dapat dibagi menjadi lima bagian
yaitu:
\begin{enumerate}[label=\arabic*.]
\item Penyusunan vektor beban dan matriks kekakuan elemen.
\item Penyusunan vektor beban dan matriks kekakuan global.
\item Penyelesaian persamaan kekakuan yang menghasilkan perpindahan global.
\item Perhitungan deformasi elemen dengan menggunakan transformasi koordinat.
\item Penentuan tegangan atau gaya-gaya elemen.
\end{enumerate}

Kunci penyederhanaan proses perakitan matriks kekakuan titik $SJ$ ialah
pemakaian matriks kekakuan batang untuk aksi dan perpindahan di kedua ujung
setiap batang. Jika perpindahan batang didasarkan pada koordinat struktur
global, maka ia akan berimpit dengan perpindahan titik kumpul. Dalam hal
ini semua komplikasi geometris harus ditangani secara setempat, dan
transfer informasi batang ke array struktur bersifat langsung, yaitu
matriks kekakuan dan vektor beban ekivalen dapat dirakit dengan penjumlahan
langsung sebagai ganti perkalian matriks. Dengan demikian, perakitan
matriks kekakuan titik untuk $m$ batang dapat dituliskan sebagai:
\begin{equation}
SJ = \sum_{i=1}^{m} SMS_i
\label{eq:2-1}
\end{equation}
dalam persamaan ini, simbol $SMS_i$ menyatakan matriks kekakuan batang
ke-$i$ dengan gaya ujung dan perpindahan (untuk kedua ujung) yang dalam
arah sumbu struktur. Agar dapat dijumlahkan, semua matriks kekakuan batang
harus diperluas ke ukuran yang sama seperti $SJ$ dengan menambahkan baris
dan kolom nol. Akan tetapi, operasi ini dapat dihindari dalam pembuatan
program dengan meletakan elemen dari $SMS_i$ ke lokasi yang sesuai pada
$SJ$. Demikian juga halnya, vektor beban ekivalen $AE$ dapat dibentuk dari
kontribusi batang sebagai berikut:
\begin{equation}
AE = \sum_{i=1}^{m} AMS_i
\label{eq:2-2}
\end{equation}
dalam hal ini $AMS_i$ adalah vektor gaya (aksi) jepit ujung dalam arah
sumbu struktur di kedua ujung batang $i$. Seperti pada $SMS_i$, secara
teoritis $AMS_i$ perlu diperbesar dengan bilangan nol agar dapat dijumlahkan
secara matriks, tetapi langkah ini bisa dihindari dalam program. Beban
titik tumpul ekivalen dalam persamaan \eqref{eq:2-2} kemudian dijumlahkan
dengan beban sebenarnya yang diberikan di titik kumpul untuk membentuk
vektor beban total atau gabungan.

Untuk memisahkan suku yang berkaitan dengan perpindahan bebas struktur dari
perpindahan pengekang tumpuan (\textit{support restraint}), matriks
kekakuan dan beban harus ditata ulang (\textit{rearranged}) dan disekat
(\textit{partitioned}). Penataan ulang suku-suku ini dapat dilakukan
setelah perakitan selesai. Dalam program komputer, penataan ulang ini lebih
mudah dilakukan bila informasi ditransfer dari array batang yang kecil ke
array struktur yang besar.

Setelah matriks kekakuan dan vektor beban ditata ulang (\textit{rearranged})
dan disekat-sekat (\textit{partitioned}), maka penyelesaian dasar untuk
perpindahan pada titik kumpul yang bebas akibat beban adalah:
\begin{equation}
DF = SFF^{-1} \, AFC
\label{eq:2-3}
\end{equation}
dalam persamaan ini $AFC$ adalah vektor beban titik kumpul gabungan
(sebenarnya dan ekivalen) yang selaras (\textit{correspond}) dengan $DF$.
Walaupun secara simbolis penyelesaian dalam persamaan \eqref{eq:2-3} lebih
mudah dituliskan sebagai hasil inversi matriks kekakuan $SFF$, tetapi
perhitungan $SFF^{-1}$ dalam program komputer hakekatnya tidak efisien.
Sebagai gantinya, matriks koefisien $SFF$ difaktorisasi dan penyelesaian
untuk $DF$ diperoleh dari sapuan (\textit{sweep}) kemuka dan belakang --
dalam program optimasi beton bertulang pada struktur portal ruang ini,
faktorisasi tersebut dilakukan dengan pendekatan Choleski yang dimodifikasi
oleh fungsi \texttt{banfac()} dan penyelesaian sapuan oleh fungsi
\texttt{bansol()} pada \texttt{Solver.hpp}.

Bagian lain yang hendak dicari adalah reaksi tumpuan dan gaya jepit ujung.
Bila hanya pengaruh beban yang diperhitungkan, persamaan untuk reaksi
adalah:
\begin{equation}
AR = SRF \, DF + ARC
\label{eq:2-4}
\end{equation}
Selain itu gaya ujung batang dituliskan sebagai:
\begin{equation}
AM_i = SM_i \, DM_i + AML_i
\label{eq:2-5}
\end{equation}
Banyak variasi dalam mengorganisasi metoda kekakuan bisa dilakukan untuk
pemrograman. Tahapan analisa di atas merupakan pendekatan teratur yang
memiliki beberapa segi menguntungkan bila berhadapan dengan masalah yang
besar dan rumit (Weaver dan Gere, 1980, halaman 151).

Untuk menganalisa struktur portal ruang digunakan matrik kekakuan batang
yang terdapat pada tabel 2-1 dan dijelaskan pada bagian \ref{sec:kekakuan-batang}.

\subsection{Kekakuan Batang Portal Ruang}
\label{sec:kekakuan-batang}

Untuk membentuk matriks kekakuan titik dengan proses penjumlahan yang
ditunjukan oleh persamaan \eqref{eq:2-1}, matriks kekakuan batang yang
lengkap dalam sumbu arah struktur perlu diturunkan dahulu. Selain itu,
matriks kekakuan batang $SM_i$ untuk sumbu arah batang diperlukan dalam
menghitung gaya ujung batang dengan persamaan \eqref{eq:2-5}. Kekakuan
batang terhadap sumbu batang selalu dapat diperoleh dan ditransformasi ke
sumbu struktur.

Portal ruang merupakan kasus yang paling umum untuk struktur rangka,
batang dapat mengalami sembarang kedua belas perpindahan, dengan demikian,
matriks kekakuan batang ($SM_i$) untuk struktur portal ruang berordo
$12 \times 12$ yang dapat dilihat dalam tabel 2-1, dan setiap kolom pada
matriks menyatakan aksi akibat salah satu perpindahan satuan.

Matriks kekakuan batang portal ruang diperlihatkan pada tabel 2-1, matriks
ini disekat untuk memisahkan bagian yang berkaitan dengan kedua ujung
batang, secara umum bentuknya adalah:
\begin{equation}
SM_i = \begin{bmatrix} SM_{jj} & SM_{jk} \\ SM_{kj} & SM_{kk} \end{bmatrix}
\label{eq:2-6}
\end{equation}
pada persamaan di atas subskrip $j$ dan $k$ pada submatriks di atas
menunjukan ujung batang (Weaver dan Gere, 1980, halaman 154).

% Tabel 2-1: Matriks Kekakuan Batang Portal Ruang
%
% The original "Tabel 2-1" was an embedded raster/vector image
% (Isi/Tabel 2.doc contains only a single [pic] anchor, no extractable
% text), so it could not be recovered from the .doc binary at all. The
% matrix below is reconstructed from the standard 12x12 local-axis space
% frame member stiffness matrix as published in Weaver, W. Jr. and Gere,
% J.M., "Matrix Analysis of Framed Structures" (cited throughout this
% chapter as the method's source), consistent with the degree-of-freedom
% ordering (u,v,w,theta_x,theta_y,theta_z per end) used in Struktur.hpp /
% Variabel.hpp (SM, SMRT arrays). VERIFY AGAINST THE ORIGINAL DOCUMENT
% before relying on the exact sign convention for structural design.
\begin{table}[htbp]
\centering
\caption{Matriks Kekakuan Batang Portal Ruang (direkonstruksi dari Weaver \& Gere, 1980)}
\label{tab:2-1}
\begin{align*}
a&=\frac{EA}{L}, &
t&=\frac{GJ}{L}\\
b_z&=\frac{12EI_z}{L^3}, & c_z&=\frac{6EI_z}{L^2}, & d_z&=\frac{4EI_z}{L}, & e_z&=\frac{2EI_z}{L}\\
b_y&=\frac{12EI_y}{L^3}, & c_y&=\frac{6EI_y}{L^2}, & d_y&=\frac{4EI_y}{L}, & e_y&=\frac{2EI_y}{L}
\end{align*}
\resizebox{\textwidth}{!}{$
SM_i =
\begin{array}{c|cccccccccccc}
 & u_j & v_j & w_j & \theta_{xj} & \theta_{yj} & \theta_{zj} & u_k & v_k & w_k & \theta_{xk} & \theta_{yk} & \theta_{zk} \\ \hline
u_j          &  a &    0 &    0 &  0 &    0 &    0 & -a &    0 &    0 &  0 &    0 &    0 \\
v_j          &  0 &  b_z &    0 &  0 &    0 &  c_z &  0 & -b_z &    0 &  0 &    0 &  c_z \\
w_j          &  0 &    0 &  b_y &  0 & -c_y &    0 &  0 &    0 & -b_y &  0 & -c_y &    0 \\
\theta_{xj}  &  0 &    0 &    0 &  t &    0 &    0 &  0 &    0 &    0 & -t &    0 &    0 \\
\theta_{yj}  &  0 &    0 & -c_y &  0 &  d_y &    0 &  0 &    0 &  c_y &  0 &  e_y &    0 \\
\theta_{zj}  &  0 &  c_z &    0 &  0 &    0 &  d_z &  0 & -c_z &    0 &  0 &    0 &  e_z \\
u_k          & -a &    0 &    0 &  0 &    0 &    0 &  a &    0 &    0 &  0 &    0 &    0 \\
v_k          &  0 & -b_z &    0 &  0 &    0 & -c_z &  0 &  b_z &    0 &  0 &    0 & -c_z \\
w_k          &  0 &    0 & -b_y &  0 &  c_y &    0 &  0 &    0 &  b_y &  0 &  c_y &    0 \\
\theta_{xk}  &  0 &    0 &    0 & -t &    0 &    0 &  0 &    0 &    0 &  t &    0 &    0 \\
\theta_{yk}  &  0 &    0 & -c_y &  0 &  e_y &    0 &  0 &    0 &  c_y &  0 &  d_y &    0 \\
\theta_{zk}  &  0 &  c_z &    0 &  0 &    0 &  e_z &  0 & -c_z &    0 &  0 &    0 &  d_z
\end{array}
$}
\end{table}


\section{Struktur Beton Bertulang}

\subsection{Tinjauan Umum}

Beton bertulang telah digunakan sebagai material bangunan di setiap negara.
Pada banyak negara seperti Amerika dan Kanada, beton bertulang merupakan
material struktur yang dominan dalam konstruksi bangunan. Keuntungan dari
penggunaan beton bertulang adalah tulangan baja dan bahan pembentuk beton
tersedia di banyak tempat, hanya dibutuhkan kemampuan yang relatif rendah
dalam konstruksi beton, dan dari segi ekonomi beton bertulang lebih
menguntungkan dibandingkan bahan konstruksi lainnya (MacGregor, 1997,
halaman 1).

Pada struktur beton bertulang, balok utama yang langsung ditumpu oleh
kolom dianggap menyatu secara kaku dengan kolom. Sistem kolom dan balok
induk seperti ini dikatakan sebagai sistem portal. Sistem portal telah
lama digunakan sebagai sistem bangunan yang dapat menahan beban vertikal
gravitasi dan lateral akibat gempa. Sistem ini memanfaatkan kekakuan
balok-balok utama dan kolom. Dengan demikian integritas antara balok utama
dan kolom harus mendapat perhatian dan pendetailan tersendiri, karena di
sekitar daerah ini timbul gaya geser dan momen yang besar yang dapat
menyebabkan retak atau patahnya penampang (Wahyudi dan Rahim, 1997,
halaman 14).

\subsection{Balok}

Balok adalah anggota struktur yang paling utama mendukung beban luar serta
berat sendirinya oleh momen dan gaya geser (MacGregor, 1997, halaman 85).

Dua hal utama yang dialami oleh suatu balok adalah kondisi tekan dan
tarik, yang antara lain karena adanya pengaruh lentur ataupun gaya
lateral. Padahal dari data percobaan diketahui bahwa kuat tarik beton
sangatlah kecil, kira-kira 10\%, dibandingkan kekuatan tekannya. Bahkan
dalam problema lentur, kuat tarik ini sering tidak diperhitungkan, sehingga
timbul usaha untuk memasang baja tulangan pada bagian tarik guna mengatasi
kelemahan beton tersebut, menghasilkan beton bertulang (Wahyudi dan Rahim,
1997, halaman 39).

Apabila penampang tersebut dikehendaki untuk menopang beban yang lebih
besar dari kapasitasnya, sedangkan di lain pihak seringkali pertimbangan
teknis pelaksanaan dan arsitektural membatasi dimensi balok, maka
diperlukan usaha-usaha lain untuk memperbesar kuat momen penampang balok
yang sudah tertentu dimensinya tersebut. Apabila hal demikian yang
dihadapi, SK-SNI T-15-1991-03 pasal 3.3.3 ayat 4 memperbolehkan penambahan
tulangan baja tarik lebih dari batas nilai $\rho_{maks}$ bersamaan dengan
penambahan tulangan baja di daerah tekan penampang balok. Hasilnya adalah
balok dengan penulangan rangkap dengan tulangan baja tarik dipasang di
daerah tarik dan tulangan tekan di daerah tekan. Pada keadaan demikian
berarti tulangan baja tekan bermanfaat untuk memperbesar kekuatan balok
(Dipohusodo, 1994, halaman 85).

Dalam praktek, sistem tulangan tunggal hampir tidak pernah dimanfaatkan
untuk balok, karena pemasangan batang tulangan tambahan di daerah tekan,
misalnya di tepi atas penampang tengah lapangan, akan mempermudah
pengaitan sengkang (\textit{stirrup}) (Wahyudi dan Rahim, 1997, halaman
60).

\gambarhilang{Penampang balok bertulangan rangkap: tulangan tekan pada
jarak $d'$ dari serat tekan terluar, tulangan tarik pada jarak (tinggi
efektif) $d$.}{Penampang Tulangan Rangkap}

\gambarhilang{Diagram regangan penampang: regangan serat tekan beton
$\varepsilon_c = 0{,}003$; garis netral pada jarak $c$ dari serat tekan
terluar.}{Diagram Regangan}

Pada gambar 2-1 tampak suatu penampang dengan tulangan tekan yang
ditempatkan dalam jarak $d'$ dari serat tekan terluar dan tulangan tarik
dalam jarak tinggi efektif $d$. Diagram regangan yang terjadi dilukiskan
dalam gambar 2-2, dengan regangan serat tekan beton dianggap telah mencapai
regangan maksimum 0,003. Garis netral terletak pada jarak $c$ yang belum
diketahui nilainya. Melalui perbandingan segitiga dalam diagram ini, dapat
ditentukan besar regangan masing-masing tulangan, yaitu:
\begin{align}
\varepsilon_s' &= 0{,}003\,\frac{c-d'}{c} \label{eq:2-7}\\
\varepsilon_s &= 0{,}003\,\frac{d-c}{c} \label{eq:2-8}
\end{align}
pada persamaan di atas $a = \beta_1 c$, dengan $a$ tinggi blok tegangan
persegi ekivalen dan $\beta_1$ faktor reduksi tinggi blok tegangan (SK SNI
T-15-1991-03 pasal 3.3.2 butir 7.(3): $\beta_1 = 0{,}85$ untuk $f_c' \le 30$
MPa, berkurang $0{,}008$ untuk setiap kenaikan 1 MPa kuat beton, dengan
batas bawah $\beta_1 \ge 0{,}65$ -- lihat \texttt{Balok::analisa()} pada
\texttt{Balok.hpp}).

Pertama-tama, kedua tulangan $AS$ dan $AS'$ diasumsikan telah mencapai
tegangan leleh $f_y$, sehingga keseimbangan gaya horisontal
$0{,}85f_c'ab + AS'f_y = AS f_y$ memberikan tinggi blok tegangan awal:
\begin{equation}
a = \frac{(AS-AS')f_y}{0{,}85\,f_c'\,b}
\label{eq:2-9}
\end{equation}
dan substitusi $a=\beta_1 c$ ke dalam persamaan \eqref{eq:2-7} dan
\eqref{eq:2-8} memberikan regangan pada trial pertama ini:
\begin{align}
\varepsilon_s' &= 0{,}003\,\frac{a-\beta_1 d'}{a} \label{eq:2-10}\\
\varepsilon_s &= 0{,}003\,\frac{\beta_1 d - a}{a} \label{eq:2-10b}
\end{align}
Bila kedua kondisi persamaan regangan tersebut terpenuhi ($\varepsilon_s
\ge \varepsilon_y$ dan $\varepsilon_s' \ge \varepsilon_y$, dengan
$\varepsilon_y = f_y/E_s$), tegangan pada baja tulangan menjadi
$f_s = f_s' = f_y$ dengan $f_s$ adalah tegangan pada baja tulangan tarik,
$f_s'$ adalah tegangan pada baja tulangan tekan. Diagram tegangan
internalnya dapat dianggap menjadi tiga bagian yaitu:
\begin{align}
\text{pada daerah tekan beton:} \quad & CC = 0{,}85\,f_c'\,a\,b \label{eq:2-11}\\
\text{pada daerah tekan baja tulangan:} \quad & CS = AS'\,f_y \label{eq:2-12}\\
\text{pada daerah tarik baja tulangan:} \quad & T = AS\,f_y \label{eq:2-13}
\end{align}
dengan $AS'$ dan $AS$ berturut-turut menyatakan luas baja tulangan tekan
dan tarik. Berdasarkan keseimbangan horisontal gaya internal $CC+CS=T$,
dihasilkan persamaan:
\begin{equation}
0{,}85\,f_c'\,a\,b + AS'\,f_y = AS\,f_y
\label{eq:2-14}
\end{equation}
dengan tinggi blok tegangan adalah persamaan \eqref{eq:2-9}, dan momen
tahanan nominal penampang adalah:
\begin{equation}
M_n = 0{,}85\,f_c'\,a\,b\left(d-\frac{a}{2}\right) + AS'\,f_y\,(d-d')
\label{eq:2-16}
\end{equation}
Momen desain diperoleh dengan memberikan faktor reduksi $\phi$ tertentu,
yang disarankan dalam SK SNI T-15-1991-03, ke dalam persamaan
\eqref{eq:2-16} (Wahyudi dan Rahim, 1997, halaman 61--63). Standar SK SNI
T-15-1991-03 pasal 2.2.3 ayat 2 memberikan faktor reduksi kekuatan $\phi$
untuk berbagai mekanisme, antara lain:
\begin{enumerate}[label=\alph*.]
\item Lentur tanpa beban aksial \dotfill $\phi = 0{,}80$
\item Geser dan puntir \dotfill $\phi = 0{,}60$
\item Tarik aksial, tanpa dan dengan lentur (sengkang) \dotfill $\phi = 0{,}65$
\item Tekan aksial, tanpa dan dengan lentur (spiral) \dotfill $\phi = 0{,}7$
\item Tumpuan pada beton \dotfill $\phi = 0{,}70$
\end{enumerate}
Dengan demikian dapat dinyatakan bahwa kuat momen $M_R$ (kapasitas momen)
sama dengan kuat momen ideal $M_n$ dikalikan faktor $\phi$ (Dipohusodo,
1994, halaman 41) -- inilah yang dihitung sebagai variabel \texttt{FMU}
pada \texttt{Balok::analisa()}.

Bila persamaan regangan \eqref{eq:2-7} dan \eqref{eq:2-8} tidak dipenuhi,
mungkin tegangan tulangan tekan ataupun tulangan tarik tidak mencapai
tegangan leleh materialnya, sehingga dalam hal ini persamaan tersebut
tidak berlaku lagi. Untuk dapat menentukan momen ketahanan penampang,
regangan aktual harus dihitung terlebih dahulu, dengan tegangan baja
sebagai fungsi regangan dikalikan modulus elastisitas $E_s = 2 \times
10^5$ MPa:
\begin{align}
f_s' &= E_s\,\varepsilon_s' = 600\,\frac{a-\beta_1 d'}{a} \label{eq:2-17a}\\
f_s &= E_s\,\varepsilon_s = 600\,\frac{\beta_1 d - a}{a} \label{eq:2-18a}
\end{align}
(dengan $f_s', f_s$ dibatasi pada rentang $-f_y \le f_s',f_s \le f_y$).
Dengan mensubtitusi persamaan \eqref{eq:2-17a} dan \eqref{eq:2-18a} ke
dalam keseimbangan gaya horisontal $0{,}85f_c'ab + AS'f_s' = AS f_s$
(bentuk umum persamaan \eqref{eq:2-14} untuk tulangan yang belum tentu
leleh), diperoleh persamaan kuadrat dalam $a$:
\begin{equation}
\bigl(0{,}85\,f_c'\,b\bigr)\,a^2 \;+\; \bigl(600\,AS' - AS\,f_y\bigr)\,a
\;-\; 600\,\beta_1\,AS'\,d' \;=\; 0
\label{eq:2-19}
\end{equation}
diselesaikan dengan rumus abc, diambil akar yang bernilai positif (fungsi
\texttt{balok::analisa()} pada \texttt{Balok.hpp} menghitung persamaan ini
lewat variabel bantu \texttt{ASOL,BSOL,CSOL,DSOL}). Sehingga, momen
tahanan penampang menjadi:
\begin{equation}
M_n = 0{,}85\,f_c'\,a\,b\left(d-\frac{a}{2}\right) + AS'\,f_s'\,(d-d')
\label{eq:2-20}
\end{equation}
dengan $a$ hasil penyelesaian persamaan \eqref{eq:2-19} dan $f_s'$ dari
persamaan \eqref{eq:2-17a}.

Momen desain dapat diperoleh dengan memberikan faktor reduksi tertentu,
yang disarankan dalam SK SNI T-15-1991-03, ke dalam persamaan
\eqref{eq:2-20} (Wahyudi dan Rahim, 1997, halaman 61--63).

Seperti pada tulangan tunggal, keruntuhan tarik atau tekan dapat pula
terjadi pada penampang tulangan rangkap. Bila hal tersebut terjadi,
keruntuhan tarik dengan pelelehan tulangan lebih disukai daripada
keruntuhan tekan dengan kehancuran beton yang mendadak. Keadaan ini dapat
dikendalikan dengan memberikan batasan tertentu pada tulangan terpasang.
SK SNI T-15-1991-03 pasal 3.3.3 menetapkan batasan:
\begin{equation}
\rho - \rho' \le 0{,}75\,\rho_b
\label{eq:2-21}
\end{equation}
dengan $\rho_b$ adalah rasio imbang (\textit{balance ratio}) baja tulangan
yang bersesuaian dengan penampang tulangan tunggal (Wahyudi dan Rahim,
1997, halaman 64). Untuk menentukan rasio penulangan keadaan seimbang
($\rho_b$) digunakan persamaan:
\begin{equation}
\rho_b = 0{,}85\,\beta_1\,\frac{f_c'}{f_y}\,\frac{600}{600+f_y}
\label{eq:2-22}
\end{equation}
dengan $f_c'$ dan $f_y$ dalam MPa, $\beta_1$ adalah konstanta yang
merupakan fungsi dari kelas kuat beton (Dipohusodo, 1994, halaman 37).
Standar SK SNI T-15-1991-03 menetapkan nilai $\beta_1$ diambil $0{,}85$
untuk $f_c' \le 30$ MPa, berkurang $0{,}008$ untuk setiap kenaikan 1 MPa
kuat beton, dan nilai tersebut tidak boleh kurang dari $0{,}65$
(Dipohusodo, 1994, halaman 31). Program mengimplementasikan persamaan
\eqref{eq:2-21}--\eqref{eq:2-22} sebagai \texttt{RHB} dan
\texttt{kendala\_rho\_b} pada \texttt{Balok.hpp}.

\subsection{Kolom Biaksial}

Kolom adalah batang tekan vertikal dari rangka (\textit{frame}) struktur
yang memikul beban dari balok. Kolom meneruskan beban-beban dari elevasi
atas ke elevasi yang lebih bawah hingga akhirnya sampai ke tanah melalui
fondasi. Karena kolom merupakan komponen tekan, maka keruntuhan pada suatu
kolom merupakan lokasi kritis yang dapat menyebabkan keruntuhan
(\textit{collapse}) lantai yang bersangkutan, dan juga merupakan batas
runtuh total (\textit{ultimate total collapse}) seluruh strukturnya. Oleh
karena itu dalam merencanakan kolom perlu lebih teliti, yaitu dengan
memberikan kekuatan cadangan yang lebih tinggi daripada yang dilakukan
pada balok dan elemen struktur horisontal lainnya, terlebih lagi
keruntuhan tekan tidak memberikan peringatan awal yang cukup jelas (Nawy,
1990).

SK SNI T-15-1991-03 (PU 1991) memberikan definisi, kolom adalah komponen
struktur bangunan yang tugas utamanya menyangga beban tekan aksial dengan
bagian tinggi yang tidak ditopang paling tidak tiga kali dimensi lateral
terkecil. Sedangkan komponen struktur yang menahan beban aksial vertikal
dengan rasio bagian tinggi dengan dimensi terkecil kurang dari tiga
disebut pedestal (Dipohusodo, 1994, halaman 287).

Bila eksentrisitas beban mempunyai harga kecil sehingga gaya aksial tekan
menjadi penentu, dan juga bila dikehendaki suatu kolom beton dengan
penampang lintang yang lebih kecil, maka umumnya distribusi tulangan
lebih baik dibuat merata di sekeliling sisi penampang tersebut. Untuk
distribusi tulangan semacam ini, baja tulangan yang terletak di bagian
tengah penampang akan menerima tegangan yang lebih kecil dibandingkan
tulangan lainnya. Ketika kapasitas ultimit kolom tersebut telah tercapai,
tegangan pada baja tulangan belum tentu mencapai tegangan lelehnya,
sedangkan baja tulangan yang berada di tepi kemungkinan besar sudah leleh
(Wahyudi dan Rahim, 1997, halaman 215).

Menurut Winter dan Nilson pada tahun 1994, mekanisme gaya aksial yang
bekerja bersamaan dengan lentur pada kedua arah dari sumbu utama penampang
terjadi pada kolom-kolom sudut bangunan, pada balok-balok yang membentuk
rangka dengan kolom dan menyalurkan momen-momen ujungnya ke kolom dalam
dua bidang yang tegak lurus. Keadaan yang sama juga dapat terjadi pada
kolom-kolom sebelah dalam, khususnya pada tata letak kolom yang tidak
teratur.

Kolom-kolom seperti pada sudut bangunan pada struktur rangka memikul gaya
aksial dan sekaligus momen dari dua sumbu aksis yang disebut kolom
biaksial. Untuk menyelesaikan masalah kolom biaksial pada kolom persegi,
dapat digunakan prosedur yang umum dipakai yaitu eksentrisitas biaksial,
$e_x$ dan $e_y$, digantikan dengan eksentrisitas uniaksial ekivalen
$e_{0x}$ atau $e_{0y}$, dan kolom didesain sebagai kolom uniaksial
(MacGregor, 1997, halaman 466).

Apabila eksentrisitas pada arah $x$ ($e_x$) dibagi dengan sisi pada arah
$x$ lebih besar daripada eksentrisitas pada arah $y$ ($e_y$) dibagi sisi
pada arah $y$, maka momen desain efektif dapat dihitung dengan:
\begin{equation}
M_{ox} = M_x + M_y\,\frac{1-\beta}{\beta}
\label{eq:2-23}
\end{equation}
Dengan cara yang sama apabila eksentrisitas pada arah $y$ dibagi dengan
sisi pada arah $y$ lebih besar daripada eksentrisitas arah $x$ dibagi sisi
arah $x$, maka momen desain efektifnya adalah:
\begin{equation}
M_{oy} = M_y + M_x\,\frac{1-\beta}{\beta}
\label{eq:2-24}
\end{equation}
dengan $M_x$ adalah momen yang terjadi pada arah $x$ dan $M_y$ adalah
momen pada arah $y$, $b$ adalah sisi efektif pada arah $y$, $h$ adalah
sisi efektif pada arah $x$, sedangkan $\beta$ adalah koefisien biaksial
yang untuk kolom persegi dengan tulangan simetris dapat dihitung dengan
rumus empiris (Hulse dan Mosley, 1986, halaman 163):
\begin{equation}
\beta = 0{,}3 + \frac{0{,}7}{0{,}6}\left(0{,}6 - \frac{P_u}{h\,b\,f_c'}\right)
\label{eq:2-25}
\end{equation}
dengan $\beta$ tidak boleh lebih kecil daripada $0{,}3$, dan $P_u$ adalah
gaya aksial. Persamaan \eqref{eq:2-23}--\eqref{eq:2-25} diimplementasikan
sebagai \texttt{MOX}, \texttt{MOY}, dan \texttt{beta} pada
\texttt{kolom::analisa()} di \texttt{Kolom.hpp}.

\section{Metoda Optimasi Struktur}

\subsection{Tinjauan Umum}

Penggunaan metoda optimasi dalam perencanaan struktur sebenarnya bukanlah
merupakan hal yang baru dan sudah banyak dikembangkan karena manfaatnya
yang banyak dirasakan. Pada tahun 1890 Maxwell mengemukakan beberapa teori
tentang desain yang rasional pada suatu struktur yang kemudian dikembangkan
lebih lanjut oleh Michell pada tahun 1904 (Wu, 1986, halaman 1). Beberapa
penelitian tentang optimasi struktur yang ditujukan untuk penggunaan
praktis telah dilakukan sekitar tahun 1940 dan 1950. Pada tahun 1960
Schmit mendemonstrasikan penggunaan teknik pemrograman non-linier untuk
desain struktur dan menyebutnya dengan istilah ``sintesa struktur'' (Wu,
1986, halaman 1). Komputer digital yang kemudian dibuat dan mampu untuk
memecahkan masalah numeris dalam skala besar telah memberikan momentum
yang besar untuk penelitian. Pada awal tahun 1970 optimasi struktur telah
menjadi sesuatu yang penting dalam berbagai aspek perancangan suatu
struktur (Wu, 1986, halaman 1).

Ada dua pendekatan utama dalam optimasi struktur. Pendekatan yang satu
menggunakan pemrograman matematika dan pendekatan yang lain menggunakan
metoda kriteria optimal. Kedua pendekatan ini masing-masing mempunyai
kelebihan dan kekurangan.

Setiap struktur rangka memiliki empat hal pokok dengan empat hal ini
dapat merupakan suatu variabel yang dapat diubah-ubah untuk mengoptimasi
struktur tersebut (variabel desain). Empat hal ini adalah ukuran elemen,
geometri struktur (posisi titik-titik kumpul), gambaran struktur
(bagaimana titik-titik kumpul tersebut dihubungkan oleh elemen-elemen),
dan material bangunan. Material bangunan biasanya ditentukan terlebih
dahulu. Persyaratan-persyaratan yang harus dipenuhi oleh struktur disebut
kendala. Dalam sebagian besar kasus, kendala berhubungan dengan kekuatan
dan defleksi struktur.

Dalam banyak kasus optimasi struktur yang telah dipecahkan sejauh ini,
variabel desain diasumsikan dapat berubah secara kontinyu. Akan tetapi
dalam desain yang sesungguhnya kadang-kadang dijumpai variabel desain
diskrit yang terbatas. Balok baja dan balok beton bertulang hanya ada
dalam ukuran standar. Untuk mengatasi masalah diskrit ini, beberapa usaha
telah dilakukan untuk memecahkan kasus-kasus tertentu, antara lain
menggunakan pendekatan \textit{branch-and-bound} (Wu, 1986, halaman 4).

Keberhasilan optimasi struktur sejauh ini masih terbatas. Pada satu sisi,
optimasi telah berhasil digunakan dalam perancangan berbagai jenis
struktur seperti \textit{overhead cranes}, bangunan-bangunan standard,
menara transmisi, girder dengan bentang pendek dan medium, dan bermacam-
macam kendaraan termasuk pesawat terbang, mobil, dan kapal. Pada sisi
lain, teknik-teknik optimasi masih kurang dapat diterapkan jika struktur
sangat besar atau jika kendala-kendalanya terlalu kompleks, belum lagi
jika ada pertimbangan-pertimbangan lain seperti standar produksi,
pengaruh estetika struktur, dan praktek-praktek konvensional dalam
industri (Wu, 1986, halaman 5).

Kesulitan utama dari penggunaan praktis optimasi struktur adalah lamanya
proses komputasi sehingga biaya perhitungan menjadi relatif tinggi. Akan
tetapi belakangan ini biaya perhitungan telah turun secara drastis karena
komputer pribadi makin cepat dan makin murah harganya. Dengan demikian
optimasi struktur yang sangat tergantung pada komputer akan semakin luas
aplikasinya.

Formulasi masalah untuk optimasi struktur dengan variabel desain yang
dapat berubah secara kontinyu dan menggunakan analisa struktur elastik
adalah program non-linier. Pendekatan untuk pemecahannya menggunakan
program linier sekuensial dan program non-linier (Wu, 1986, halaman 6).

Pendekatan linier sekuensialnya adalah melinierkan fungsi kendala
non-linier dan fungsi sasaran kemudian menyelesaikan masalah menggunakan
program linier berulang-ulang. Romstad dan Wang pada tahun 1967 dan 1968
memberikan penjelasan fisik tentang formulasi ini yaitu: yang tidak
diketahui adalah penambahan variabel desain dari suatu titik dalam daerah
layak yang diperoleh dari iterasi sebelumnya, dan kendala dari pendekatan
linier ini adalah perpindahan ijin dari titik kumpul dan gaya-gaya yang
terjadi pada elemen tidak boleh melebihi yang diijinkan. Pendekatan
non-liniernya adalah menyelesaikan masalah optimasi dengan menggunakan
program non-linier secara langsung. Akan tetapi persamaan analisa elastik
menunjukkan masalah yang besar sebab tidak ada algoritma pada program
non-linier yang dapat menangani kendala persamaan secara efisien. Teknik
yang biasa digunakan untuk menangani masalah ini adalah menggunakan
matriks untuk analisa struktur dalam algoritmanya (Wu, 1986, halaman 7).

Formulasi masalah untuk optimasi struktur dengan variabel desain yang
dapat berubah secara kontinyu dan menggunakan analisa struktur plastis
(analisa batas) adalah program linier (Wu, 1986, halaman 8). Masalah ini
dapat dipecahkan dengan mudah bahkan untuk struktur yang besar. Akan
tetapi analisa batas belum dipakai secara luas dalam desain struktur. Di
samping itu banyak kendala lain seperti defleksi tidak dapat dipenuhi
jika perancangan didasarkan pada analisa batas.

Optimasi struktur dengan sebagian atau seluruh variabel desain merupakan
nilai diskrit yang terbatas dan menggunakan analisa struktur elastis
memberikan formulasi masalahnya berupa masalah non-linier diskrit.
Beberapa pendekatan untuk memecahkan masalah ini telah diusulkan termasuk
program integer sekuensial, algoritma penelusuran arah variabel diskrit,
dan lain-lain (Wu, 1986, halaman 8).

Program integer sekuensial adalah ekstrapolasi dari program linier
sekuensial ke program diskrit. Dalam setiap iterasi, program linier
integer digunakan untuk menjamin penelusuran hanya berkisar dari satu
titik diskrit ke yang lain. Toakley pada tahun 1968 menemukan bahwa
pendekatan ini tidak efisien dan hasilnya tidak bisa diramalkan.
Reinschmidt pada tahun 1971 juga mempunyai kesimpulan yang sama meskipun
dia telah menunjukkan beberapa contoh kasus numerik yang dipecahkan
dengan cara ini, termasuk desain elastis rangka batang yang terdiri atas
9 elemen. Saglam pada tahun 1976 menggunakan pendekatan yang mirip dan
telah memecahkan beberapa contoh rangka batang.

Algoritma penelusuran arah variabel diskrit dikemukakan oleh Lai dan
Achenbach (1973, halaman 119--131). Struktur portal, kantilever, dan
pelat lantai dua lapis dengan kendala dinamis telah dioptimasi dengan
metoda ini.

Optimasi struktur dengan sebagian atau seluruh variabel desain merupakan
harga diskrit yang terbatas dan menggunakan analisa struktur plastis,
formulasi masalahnya adalah program integer atau program integer
campuran. Toakley (1968, halaman 1219) merumuskan masalah desain plastis
dengan variabel diskrit sebagai program integer dan telah dipecahkan. Fu
dan You pada tahun 1976 menggunakan metoda complex, tetapi mereka
menemukan metoda ini konvergen prematur, kadang-kadang jauh dari nilai
optimum. Levey dan Fu (1979, halaman 363--368) menggunakan metoda
complex-simplex. Lev (1977, halaman 365--371) mengusulkan algoritma
\textit{branch-and-bound} untuk memecahkan masalah diskrit dengan kendala
linier. Algoritma ini mengatasi keterbatasan dari beberapa algoritma
program integer yang lain.

Selama perkembangan optimasi struktur, para ilmuwan secara bertahap
mengembangkan banyak cara yang lain untuk merumuskan dan memecahkan
berbagai masalah sehingga banyak algoritma-algoritma yang diusulkan.
Sementara situasi ini dikenal secara umum sebagai sesuatu yang tidak dapat
dihindari dan harus dialami selama tahap awal perkembangan bidang
teknologi, ada juga saran untuk optimasi struktur dengan pendekatan yang
disatukan dan sistematis.

Dalam pendekatan sistematis yang digambarkan oleh Morris dan kawan-kawan
pada tahun 1982, algoritma dibagi dalam tiga bagian, tiap bagian mempunyai
fungsi yang terdefinisi secara baik, dan program komputer untuk tiap
bagian dapat dikembangkan dan diuji secara terpisah. Pendekatan ini
memperbolehkan penggunaan program secara berdiri sendiri yang kemudian
dapat dikumpulkan dengan cepat untuk persoalan desain yang spesifik.
Pendekatan ini dapat juga digunakan sebagai bahan perbandingan untuk
perkembangan langkah-langkah optimasi yang baru. Di samping itu,
pendekatan sistematik ini dapat dengan mudah diterapkan pada masalah yang
lebih kompleks di bidang teknik yang lain (Wu, 1986, halaman 10).

Menurut Kirsch pada tahun 1981 biasanya dalam suatu perencanaan terdiri
atas empat langkah yaitu:
\begin{enumerate}[label=\arabic*.]
\item Perumusan syarat-syarat fungsional, yaitu mencari dan merumuskan
      syarat-syarat fungsional yang dalam beberapa kasus tidak terlihat
      secara nyata.
\item Perencanaan dasar, misalnya pemilihan topologi, tipe struktur dan
      material.
\item Proses optimasi, yaitu untuk memperoleh kemungkinan perencanaan
      terbaik dengan kriteria, pertimbangan dan batas-batas yang ada.
\item Pendetailan, setelah seluruh penyajian optimasi, hasil yang didapat
      harus diperiksa dan dimodifikasi jika perlu.
\end{enumerate}

Berdasarkan berbagai kemajuan ilmu dan teknologi, perancangan struktur
bangunan harus direncanakan secara optimal yaitu struktur yang paling
ekonomis serta memenuhi segala persyaratan yang diinginkan. Oleh karena
itu perlu dikembangkan suatu sistem yang mampu menangani berbagai masalah
optimasi. Secara umum masalah optimasi ada empat jenis, yaitu:
\begin{enumerate}[label=\arabic*.]
\item Optimasi bentuk.
\item Optimasi topologi.
\item Optimasi geometri.
\item Optimasi ukuran penampang.
\end{enumerate}

Dalam metoda optimasi terdapat tiga besaran utama, yaitu:
\begin{enumerate}[label=\arabic*.]
\item \textbf{Variabel desain.} Besaran yang tidak berubah nilainya
      disebut parameter tetap, sedangkan yang nilai berubah selama proses
      optimasi disebut variabel desain. Variabel desain merupakan
      variabel yang dicari dalam masalah optimasi. Contohnya adalah
      ukuran komponen struktur dan geometri struktur. Data variabel
      desain ada dua macam, yaitu data diskrit dan data kontinu. Dalam
      beberapa kasus, khususnya optimasi bentuk dan geometri, variabel
      desain lebih sesuai dinyatakan sebagai variabel desain kontinu
      dibandingkan variabel diskrit.
\item \textbf{Fungsi kendala.} Fungsi kendala merupakan suatu fungsi yang
      memberikan batasan daerah layak dan daerah tak layak. Dalam bidang
      teknik terdapat dua macam kendala yaitu:
      \begin{enumerate}[label=\alph*.]
      \item Kendala rencana, yaitu kendala yang menentukan variabel
            desain selain yang memberikan batasan berdasarkan sifat.
            Kendala ini biasanya dapat dilihat secara nyata, misalnya
            batasan karena masalah fungsional, fabrikasi atau keindahan.
            Contoh kendala rencana adalah ketebalan plat, kemiringan
            atap.
      \item Kendala sifat, yaitu kendala yang didapat dari persyaratan
            sifat. Biasanya kendala ini tidak dapat terlihat secara nyata
            karena berhubungan dengan analisis struktur. Contoh kendala
            sifat adalah batas tegangan maksimum, perpindahan
            (\textit{displacements}) yang diijinkan, kekuatan tekuk.
      \end{enumerate}
\item \textbf{Fungsi sasaran.} Fungsi sasaran adalah suatu fungsi yang
      mengandung kriteria dari struktur yang diinginkan, misalnya
      struktur dengan berat paling ringan, dengan harga termurah, paling
      aman atau paling efisien. Pemilihan fungsi sasaran merupakan hal
      yang terpenting dalam proses optimasi agar dapat mencapai sasaran
      yang sebenarnya sedekat mungkin. Dalam beberapa situasi fungsi
      sasaran dapat terlihat jelas. Misalnya jika ingin mencari harga
      yang termurah maka fungsi sasarannya dapat diasumsikan ke dalam
      berat strukturnya. Namun terkadang sulit juga untuk menentukan
      harga yang sebenarnya dari sebuah konstruksi, misalnya struktur
      dengan berat paling ringan belum tentu yang termurah, karena
      biasanya masalah harga minimum akan termasuk juga harga bahan,
      fabrikasi, transportasi dan lain-lain. Masalah dalam optimasi dapat
      dibedakan menjadi dua, yaitu:
      \begin{enumerate}[label=\alph*.]
      \item Masalah optimasi linier merupakan dasar dari metoda optimasi
            secara matematis. Dalam masalah ini fungsi kendala dan fungsi
            sasarannya semuanya dinyatakan dalam fungsi linier. Fungsi
            kendala dapat berupa persamaan maupun pertidaksamaan, dan
            fungsi sasarannya berupa meminimumkan dan memaksimumkan.
      \item Masalah optimasi tak linier, yaitu bila fungsi kendala dan
            fungsi sasarannya tak linier.
      \end{enumerate}
\end{enumerate}

Masalah optimasi dalam bidang teknik, pada umumnya berupa masalah
optimasi tak linier. Masalah yang tak linier ini juga dapat dilinierkan,
tetapi memberikan hasil yang kurang akurat ditinjau dari segi teknik.
Oleh karena itu terpaksa diselesaikan memakai program tak linier yang
lebih sulit dipelajari dibandingkan program linier, karena memerlukan
matematika yang kompleks.

Penyelesaian masalah optimasi dapat dipakai dua cara yaitu:
\begin{enumerate}[label=\arabic*.]
\item \textbf{Metoda analisa.} Metoda ini menggunakan dasar teori
      matematika yang dibuat oleh Maxwell pada tahun 1890 dan Michell
      tahun pada 1904 dan memberikan hasil eksak namun hanya dapat
      digunakan untuk masalah optimasi yang sederhana saja karena pada
      beberapa masalah yang lebih kompleks pengolahan matematikanya
      sangat tidak sederhana (Wibowo, 1996, halaman 2).
\item \textbf{Metoda numerik.} Metoda optimasi numerik berkembang sejak
      ditemukannya komputer sebagai alat bantu hitung. \textit{Dynamic
      programming, integer programming, steepest descent, sequential
      unconstraint minimization technique, gradient projection}, dan
      \textit{penalty function} merupakan metoda optimasi numerik yang
      sering dipakai untuk menyelesaikan masalah optimasi di bidang
      sipil (Wibowo, 1996, halaman 3). Dalam metoda ini nilai yang akan
      dicari didekati dengan cara iterasi dan proses iterasi dihentikan
      apabila nilai yang dicari sudah cukup dekat dengan titik optimal
      yang sesungguhnya (Kirsch, 1981, halaman 5).
\end{enumerate}

Metoda-metoda tersebut di atas mempunyai kelemahan, yaitu mempunyai
peluang relatif besar untuk konvergen ke titik optimum lokal, bila
dipakai untuk menyelesaikan masalah optimasi dengan jumlah titik optimum
lebih dari satu. Hal ini dapat dimaklumi, karena penurunan perumusannya
berdasarkan asumsi fungsi konveks. Selain itu juga memerlukan analisis
struktur yang banyak, keharusan mengikutsertakan semua kendala dalam
model matematikanya dan modifikasi variabel yang kurang efisien (Wibowo,
1996, halaman 3).

\subsection{Metoda Optimasi Polihedron Fleksibel}
\label{sec:polihedron}

Metoda polihedron fleksibel merupakan metoda pengembangan dari metoda
\textit{simplex}, yang dikembangkan oleh Nelder-Mead (Haftka, 1991,
halaman 64), dasar pemikiran metoda \textit{simplex} adalah menurunkan
nilai fungsi sasaran secara kontinu dimulai dari suatu nilai fungsi awal
sampai mencapai nilai fungsi minimum terpenuhi (Haftka, 1991, halaman
64).

Metoda ini bermanfaat untuk mencari harga-harga ekstrem suatu fungsi
dengan banyak variabel, dengan turunan dari fungsi tersebut sulit untuk
dicari. Metoda polihedron fleksibel menggunakan pencerminan
(\textit{reflection}), ekspansi (\textit{expansion}) dan penyusutan
(\textit{contraction}) dalam melakukan penelusuran.

Polihedron Fleksibel menggunakan banyak titik coba, dengan jumlah titik
coba $N_{titik}$ minimum sama dengan jumlah variabel desain $JVD$
ditambah satu. $N_{titik} = JVD+1$ dipakai oleh Harb dan Mattews pada
tahun 1987. Box mengusulkan $N_{titik} = 2 \cdot JVD$, sedangkan
$N_{titik} = JVD+2$ diusulkan oleh Biles pada tahun 1983. Jumlah titik
coba makin banyak dimaksudkan untuk mengurangi terjadinya konvergen
prematur, tetapi memperlambat konvergensi.

Masalah meminimumkan fungsi sasaran $f=f(x,y)$, dengan dua variabel
desain $x$ dan $y$, dan memakai jumlah titik coba $N_{titik} = JVD+1 = 3$,
prosedurnya adalah sebagai berikut:
\begin{enumerate}[label=\arabic*.]
\item Tentukan atau acak tiga buah titik di dalam ruang variabel disain,
      kemudian hitung nilai fitnessnya. Titik terbaik disebut titik $B$
      (\textit{best}) dengan koordinat $(X_B,Y_B)$, titik terjelek
      disebut titik $W$ (\textit{worst}) dengan koordinat $(X_W,Y_W)$,
      sedang titik lainnya disebut titik $G$ (\textit{good}) dengan
      koordinat $(X_G,Y_G)$.
\item Hitung titik pusat $M$ dengan koordinat $(X_M,Y_M)$ yang didapat
      dengan merata-rata titik coba selain titik $W$, sehingga:
      \begin{align}
      X_M &= 0{,}5\,(X_B + X_G) \label{eq:2-26}\\
      Y_M &= 0{,}5\,(Y_B + Y_G) \label{eq:2-27}
      \end{align}
\item Tentukan arah penelusuran, yaitu garis yang menghubungkan titik $W$
      dan titik $M$ dan dinyatakan sebagai:
      \begin{align}
      X_S &= X_M - X_W \label{eq:2-28}\\
      Y_S &= Y_M - Y_W \label{eq:2-29}
      \end{align}
\item Cari titik coba baru dalam arah $S$ (\textit{search}) yang
      memberikan nilai fitness lebih baik daripada fitness di titik $W$,
      titik coba ini tidak boleh identik dengan titik $M$. Titik coba
      baru $T$ (\textit{try}) dengan koordinat $X_T,Y_T$ adalah:
      \begin{align}
      X_T &= X_W + \alpha\,X_S \label{eq:2-30}\\
      Y_T &= Y_W + \alpha\,Y_S \label{eq:2-31}
      \end{align}
      $\alpha$ adalah koefisien refleksi. Kalau tidak ditemukan titik
      coba baru yang lebih baik dari titik $W$ maka dilakukan penyusutan
      menuju $B$, besarnya penyusutan sama dengan setengah jaraknya
      terhadap titik $B$, sehingga titik $W$ baru adalah:
      \begin{align}
      X_W^{\text{baru}} &= 0{,}5\,(X_W + X_B) \label{eq:2-32}\\
      Y_W^{\text{baru}} &= 0{,}5\,(Y_W + Y_B) \label{eq:2-33}
      \end{align}
      dan titik $G$ baru adalah
      \begin{align}
      X_G^{\text{baru}} &= 0{,}5\,(X_W + X_G) \label{eq:2-34}\\
      Y_G^{\text{baru}} &= 0{,}5\,(Y_W + Y_G) \label{eq:2-35}
      \end{align}
      Kemudian diperiksa apakah sudah konvergen atau belum, kalau sudah
      maka berhenti, kalau belum ulangi ke langkah dua lagi.
      Pemeriksaan konvergensi dapat menggunakan persamaan-persamaan yang
      ada di bawah ini:
      \begin{align}
      |X_B - X_G| &\le \varepsilon \label{eq:2-36a}\\
      |X_B - X_W| &\le \varepsilon \label{eq:2-36b}\\
      |X_W - X_G| &\le \varepsilon \label{eq:2-36c}\\
      |Y_B - Y_G| &\le \varepsilon \label{eq:2-36d}\\
      |Y_B - Y_W| &\le \varepsilon \label{eq:2-36e}\\
      |Y_W - Y_G| &\le \varepsilon \label{eq:2-36f}
      \end{align}
      dengan $\varepsilon$ adalah suatu nilai konvergensi dan diambil
      sekecil mungkin, misalnya $0{,}0003$.
\end{enumerate}

Secara umum metoda ini membuat sebuah segi banyak dalam ruang variabelnya
yang terus diiterasi sehingga segi banyak itu makin lama makin mengecil,
dan akan didapatkan hasil yang optimum begitu segi banyak itu menjadi
sangat kecil sekali, yang ditentukan nilainya sebagai suatu nilai
konvergensi.
